%!TEX TS-program = xelatex
%!TEX encoding = UTF-8 Unicode

%-------------------------------------------------------------------------------
% CONFIGURATIONS
%-------------------------------------------------------------------------------
\documentclass[12pt, a4paper]{awesome-cv}

% --- [中文支持配置 START] ---
\usepackage{xeCJK}

% 1. 正文主字体 (对应 AdobeSongStd-Light.otf)
\setCJKmainfont[
    Path = fonts/zh_CN-Adobe/,
    Extension = .otf,
    AutoFakeBold = true,
    AutoFakeSlant = true
]{AdobeSongStd-Light}

% 2. 标题/无衬线字体 (对应 AdobeHeitiStd-Regular.otf)
\setCJKsansfont[
    Path = fonts/zh_CN-Adobe/,
    Extension = .otf,
    AutoFakeBold = true,
    AutoFakeSlant = true
]{AdobeHeitiStd-Regular}

% 3. 等宽/代码字体 (对应 AdobeHeitiStd-Regular.otf)
\setCJKmonofont[
    Path = fonts/zh_CN-Adobe/,
    Extension = .otf,
    AutoFakeBold = true,
    AutoFakeSlant = true
]{AdobeHeitiStd-Regular}

% (可选) 定义楷体供特殊使用，如 \kai{这段字是楷体}
\setCJKfamilyfont{kai}[
    Path = fonts/zh_CN-Adobe/,
    Extension = .otf,
    AutoFakeBold = true
]{AdobeKaitiStd-Regular}
\newcommand{\kai}{\CJKfamily{kai}}

% 调整行距
\linespread{1.25}
% --- [中文支持配置 END] ---


% 配置页边距 (已在 cls 中定义，此处可微调)
\geometry{left=1.4cm, top=.8cm, right=1.4cm, bottom=1.8cm, footskip=.5cm}

% 指定英文字体路径
\fontdir[fonts/]

% 强调色设置
\colorlet{awesome}{awesome-red}

%-------------------------------------------------------------------------------
%	个人信息
%-------------------------------------------------------------------------------
% 您的照片文件名，请确保目录下有 profile.png
\photo[rectangle,edge,left]{profile.png}

\name{}{赵文彪}
\position{算法与竞赛协会组织部长{\cdotp}学习委员}
\address{新疆大学 (2024级)}

\mobile{131-8873-7995}
\email{geneva4869@163.com}
% \homepage{yourwebsite.com} 
% \github{https://github.com/winbeau} 

\vspace{-3mm}
% 引言
\quote{``代码如诗，简洁至上"}


%================= Header 字体大小 =================
\renewcommand*{\headerfirstnamestyle}[1]{{\fontsize{30pt}{1em}\headerfont\bfseries #1}}
\renewcommand*{\headerlastnamestyle}[1]{{\fontsize{30pt}{1em}\headerfont\bfseries #1}}
\renewcommand*{\headerpositionstyle}[1]{{\fontsize{12pt}{1em}\bodyfont\scshape\color{awesome}{#1}}}
\renewcommand*{\headeraddressstyle}[1]{{\fontsize{12pt}{1em}\headerfont\itshape\color{lighttext}{#1}}}
\renewcommand*{\headersocialstyle}[1]{{\fontsize{10pt}{1em}\headerfont\color{text}{#1}}}
\renewcommand*{\headerquotestyle}[1]{{\fontsize{11pt}{1.2em}\bodyfont\itshape\color{darktext}{#1}}}

%================= Header 间距 =================
\renewcommand{\acvHeaderAfterNameSkip}{5mm}
\renewcommand{\acvHeaderAfterPositionSkip}{1.2mm}
\renewcommand{\acvHeaderAfterAddressSkip}{0mm}
\renewcommand{\acvHeaderAfterSocialSkip}{3mm}
\renewcommand{\acvHeaderAfterQuoteSkip}{8mm}

%================= 正文字体（教育经历 / 技能 / 项目） =================
\renewcommand*{\entrytitlestyle}[1]{{\fontsize{13pt}{18pt}\bodyfont\bfseries\color{darktext}{#1}}}
\renewcommand*{\entrypositionstyle}[1]{{\fontsize{12pt}{17pt}\bodyfont\scshape\color{graytext}{#1}}}
\renewcommand*{\entrylocationstyle}[1]{{\fontsize{12pt}{17pt}\bodyfontlight\slshape\color{awesome}{#1}}}
\renewcommand*{\entrydatestyle}[1]{{\fontsize{12pt}{17pt}\bodyfontlight\slshape\color{graytext}{#1}}}
\renewcommand*{\descriptionstyle}[1]{{\fontsize{12pt}{18pt}\bodyfontlight\color{text}{#1}}}
\renewcommand*{\skilltypestyle}[1]{{\fontsize{12pt}{18pt}\bodyfont\bfseries\color{darktext}{#1}}}
\renewcommand*{\skillsetstyle}[1]{{\fontsize{12pt}{18pt}\bodyfontlight\color{text}{#1}}}


%================= section 标题（只调字体大小） =================
\makeatletter
\renewcommand*{\sectionstyle}[1]{%
  {\fontsize{18pt}{22pt}\bodyfont\bfseries #1}%
}
\makeatother

% ================= 修复 \cventry 间距过大问题 =================
\makeatletter
\renewcommand*{\cventry}[5]{%
  \vspace{1mm}
  \setlength\tabcolsep{0pt}
  \begin{tabular*}{\textwidth}{@{\extracolsep{\fill}} L{\textwidth - 7cm} R{7cm}}
    % 第一行：项目名称 + 右侧第一行（Location）
    \entrytitlestyle{#2} & \entrylocationstyle{#3} \\[-0.5mm]
    % 第二行：项目地址（左） + 右侧第二行（Date）
    \entrypositionstyle{#1} & \entrydatestyle{#4} \\[0mm]
    % 描述内容（如果有）
    \multicolumn{2}{L{\textwidth}}{\descriptionstyle{#5}}
  \end{tabular*}
  \vspace{-2mm}
}
\makeatother

% ================= 调整 \cvhonor 逗号填充 =================
\makeatletter
\renewcommand*{\cvhonor}[4]{%
  \honordatestyle{#4}%
  & \honorpositionstyle{#1} \honortitlestyle{#2}%
  & \honorlocationstyle{#3}%
  \\
}
\makeatother

% ====== 强化灰色对比度（更深，阅读更清晰） ======
\definecolor{darktext}{HTML}{3A3A3A}    % 主体文字（比原来深）
\definecolor{graytext}{HTML}{3E3E3E}    % 次级灰，比默认灰更深
\definecolor{lighttext}{HTML}{4E4E4E}   % 辅助灰，比默认淡灰更深

% ===== 荣誉奖项全部加黑、加粗 =====
\renewcommand*{\honordatestyle}[1]{%
  {\fontsize{12pt}{16pt}\bodyfont\bfseries\color{darktext} #1}%
}

\renewcommand*{\honorpositionstyle}[1]{%
  {\fontsize{12pt}{16pt}\bodyfont\bfseries\color{darktext} #1}%
}

\renewcommand*{\honortitlestyle}[1]{%
  {\fontsize{12pt}{16pt}\bodyfont\bfseries\color{darktext} #1}%
}

\renewcommand*{\honorlocationstyle}[1]{%
  {\fontsize{12pt}{16pt}\bodyfont\bfseries\color{darktext} #1}%
}




\begin{document}

% 渲染头部
\makecvheader[C]

% 渲染页脚
\makecvfooter
  {\today}
  {赵文彪~~~·~~~个人简历}
  {\thepage}


%-------------------------------------------------------------------------------
%	简历内容
%-------------------------------------------------------------------------------
\vspace{-8mm}
% --- [教育经历] ---
\cvsection{\textcolor{awesome}{教育}经历}
\vspace{-3mm}
\begin{cventries}
  \cventry
    {计算机科学与技术专业 / 本科大二} % 学位/年级
    {新疆大学} % 学校
    {中国, 新疆} % 地点
    {2024年 9月 - 至今} % 时间
    {
      \vspace{-4mm}
      \begin{cvitems}
        \item {核心成绩：\textbf{GPA 排名 2/116}} %
        \item {英语能力：CET-4 (514分) / CET-6 (431分)。} %
        \item {普通话等级：二级甲等} %
      \end{cvitems}
    }
\end{cventries}

\vspace{-2mm}
% --- [专业技能] ---
\cvsection{\textcolor{awesome}{专业}技能}
\begin{cvskills}
  \cvskill
    {编程语言} % Category
    {\textbf{C++} (代码量约5万行), Python (熟悉基础与 Torch、NumPy 库), LaTeX/Markdown} 

  \cvskill
    {开发环境} % Category
    {Linux/Ubuntu 命令行环境, Vim, Git}

  \cvskill
    {科研进度} % Category
    {具备一定论文阅读能力，正在进行顶刊论文复现} 
\end{cvskills}

% --- [项目与实践经历] ---
\cvsection{\textcolor{awesome}{项目}与实践经历}
\vspace{-3mm}
\begin{cventries}

  \cventry
    {项目地址} % Role
    {使用SOTA方法从0到1实现Transformer} % Project Name
    {Stanford-CS336 Assignment 1} % Location
    {github.com/winbeau/cs336-jupyter} % Date
    {}
  \vspace{-5mm}
  \cventry
    {项目地址} % Role
    {实验室资料上传与共享网站} % Project Name
    {https://selab.team} % Location
    {github.com/winbeau/KnoHub} % Date
    {}
  \vspace{-5mm}
  \cventry
    {项目地址} % Role
    {面向 KiTS19 肾脏肿瘤分割的医学影像分析平台} % Project Name
    {http://82.157.209.193} % Location
    {github.com/winbeau/KidneyTumorAI} % Date
    {}
  \vspace{-5mm}
  \cventry
    {项目地址} % Role
    {DeepSeek 驱动的 PDF 算法题翻译工具} % Project Name
    {https://algofluent.xju-wiki.top} % Location
    {github.com/winbeau/AlgoFluent} % Date
    {}
  \vspace{-5mm}
  \cventry
    {项目地址} % Role
    {Ubuntu 环境初始化方案} % Project Name
    {脚本仓库} % Location
    {github.com/winbeau/FluxEnv} % Date
    {}

\end{cventries}

\vspace{-8mm}

% --- [荣誉奖项] ---
\cvsection{\textcolor{awesome}{荣誉}奖项}
\begin{cvhonors}
  \cvhonor
    {一等奖 (排名 23/2837),} % Award
    {第十六届全国大学生数学竞赛 (非数学A类)} % Event
    {省部级} % Location
    {2024} % Date

  \cvhonor
    {全省第 19 名,} % Award
    {2025 ICPC 国际大学生程序设计竞赛 (新疆省赛)} % Event
    {省部级} % Location
    {2025} % Date

  \cvhonor
    {铜奖,} % Award
    {2025 ICPC 国际大学生程序设计竞赛全国邀请赛} % Event
    {国家级} % Location
    {2025} % Date

  \cvhonor
    {三等奖,} % Award
    {中国大学生计算机设计大赛} % Event
    {省部级} % Location
    {2025} % Date

  \cvhonor
    {优秀奖,} % Award
    {GPLT 中国高校计算机大赛-团体程序设计天梯赛} % Event
    {国家级} % Location
    {2025} % Date
  \cvhonor
    {二等奖,} % Award
    {第7届大学生智能技术应用大赛} % Event
    {国家级} % Location
    {2025} % Date
  \cvhonor
    {优秀奖,} % Award
    {第六届全国大学生算法设计与编程挑战赛(秋季赛)} % Event
    {国家级} % Location
    {2025} % Date
  \cvhonor
    {国家励志奖学金} % Award
    {} % Event
    {国家级} % Location
    {2025} % Date
  
\end{cvhonors}


\end{document}